\pdfoutput=1
\documentclass[11pt]{article}

\usepackage{ACL2023}
\usepackage{times}
\usepackage{latexsym}
\usepackage[T1]{fontenc}

\usepackage[utf8]{inputenc}
\usepackage{microtype}

%%%%%%%%%%%%%%%%%%% TITLE AND AUTHORS %%%%%%%%%%%%%%%%%%%

\title{Faithful RAG: Evaluating Citation Accuracy and Retrieval Robustness in Domain-Specific Scientific Literature}

\author{\normalfont
Elie Bruno | 355932 | \texttt{elie.bruno@epfl.ch} \\
Andrea Trugenberger | SCIPER | \texttt{andrea.trugenberger@epfl.ch} \\
Author 3 | SCIPER 3 | \texttt{author3@epfl.ch} \\
Author 4 | SCIPER 4 | \texttt{author4@epfl.ch} \\
Team Faithful RAG
}

%%%%%%%%%%%%%%%%%%% PROPOSAL %%%%%%%%%%%%%%%%%%%

\begin{document}
\maketitle

\section{Introduction}

Retrieval-Augmented Generation (RAG) has become the standard approach for grounding LLM outputs in external knowledge~\cite{retrievalaugmentedgeneration}, yet citation presence does not guarantee citation accuracy: models frequently cite the wrong passage, attribute claims to non-existent sources, or contradict their own retrieved context~\cite{min-etal-2023-factscore}. This problem is especially acute in scientific and policy domains, where a fabricated citation cannot be verified without reading the original paper.

We address this in the domain of \textbf{intellectual property and innovation policy}. We build a RAG pipeline that ingests academic papers via VLM-based PDF parsing and design a systematic evaluation of citation faithfulness and retrieval robustness. Our contributions are: (1)~a domain-specific gold benchmark of 50 expert-annotated question-answer pairs with ground-truth passages; (2)~an NLI-based citation verification pipeline following the ALCE framework~\cite{gao2023enabling}; (3)~an implementation of Corrective RAG~\cite{yan2024corrective} with iterative query refinement inspired by FLARE~\cite{jiang2023active}; and (4)~a multi-axis retrieval ablation comparing bi-encoders and late-interaction models.

\section{Methodology}

\paragraph{RAG Pipeline.} We ingest academic PDFs using a vision-language model (Dolphin~1.5) for layout-aware markdown extraction, apply hybrid semantic chunking (coarse ${\sim}$2000 and fine ${\sim}$300 chars), embed with FAISS, optionally rerank with a neural reranker, and generate cited answers via multiple LLMs.

\paragraph{Gold Dataset.} Each team member annotates 12--13 question-answer pairs from the corpus with ground-truth passages at chunk level, spanning policy impact, methodology, cross-country comparison, factual extraction, and multi-hop reasoning.

\paragraph{Retrieval Ablation.} We compare four retrieval approaches---three bi-encoders (OpenAI \texttt{text-embedding-3-small}, BGE-M3, E5-large) and ColBERTv2~\cite{santhanam2022colbertv2} for late interaction---across chunk granularities with and without reranking, measuring precision@k, recall@k, MRR, and hit rate.

\paragraph{Citation Verification.} We decompose generated answers into atomic claims via prompted extraction, then classify each claim--citation pair using DeBERTa-v3-large-mnli for NLI, labeling claims as \textit{supported}, \textit{not supported}, \textit{fabricated source}, or \textit{no citation}.

\paragraph{Corrective RAG.} A retrieval quality evaluator scores document relevance after initial retrieval. When confidence falls below a threshold, the system reformulates the query and re-retrieves. We ablate the confidence threshold.

\paragraph{End-to-End Evaluation.} We evaluate using RAGAS metrics (faithfulness, answer relevancy, context precision, context recall) and compare chunked RAG against a long-context baseline (128K window).

\section{Timeline and Team Roles}

\noindent\textbf{Weeks 1--2} (Apr 21--May 3): Repository setup, gold dataset annotation, proposal submission.\\
\textbf{Weeks 3--4} (May 5--18): Parallel implementation of all four evaluation components.\\
\textbf{Week 5} (May 19--24): Preliminary experiments; submit progress report and notebooks.\\
\textbf{Weeks 6--7} (May 26--Jun 7): Full experiments, ablations, final report.

\vspace{0.5em}
\noindent\textbf{Elie Bruno}: Retrieval ablation (embeddings, chunking, reranking).\\
\textbf{Andrea Trugenberger}: Citation verification (NLI pipeline, cross-LLM comparison).\\
\textbf{Member 3}: Corrective RAG implementation and threshold ablation.\\
\textbf{Member 4}: RAGAS evaluation and long-context baseline.\\
All members contribute to the gold evaluation dataset.

\newpage
\bibliography{anthology,custom}
\bibliographystyle{acl_natbib}

\end{document}
